\chapter*{Synthèse}
\addcontentsline{toc}{chapter}{Synthèse}
\markboth{Synthèse}{Synthèse}

Le programme \textsc{Junon} vise à construire des jumeaux numériques des milieux naturels, dont
les eaux souterraines. Prévoir l'évolution d'une nappe est aujourd'hui possible par apprentissage
automatique ; s'en servir pour décider ne l'est pas encore, car une prévision qu'on ne sait pas
justifier ne sera pas retenue à l'appui d'un arrêté de restriction d'usage. C'est à cette
difficulté, l'\textbf{explicabilité des prévisions}, que ce postdoctorat était consacré, au LIFAT
de l'Université de Tours, du 15 septembre 2025 au 30 juin 2026.

Le travail s'est heurté d'emblée à un obstacle matériel. La donnée hydrologique française est
publique, mais elle est éparpillée entre des services qui s'ignorent : les niveaux de nappe d'un
côté, les débits des rivières d'un autre, le climat ailleurs encore, sans lien entre eux et sans
indicateurs comparables d'une station à l'autre. Or on ne peut pas expliquer un modèle qu'on n'a
pas pu entraîner. \textbf{Construire ce jeu de données est donc devenu le premier travail.} Les
échanges menés avec le BRGM au début de l'année 2026 ont d'ailleurs montré que ce manque n'était
pas propre à ces travaux : le besoin d'un entrepôt consolidé et d'outils d'analyse était partagé
par les équipes du BRGM comme par les autres projets du programme.

Deux réalisations en découlent.

\paragraph{Un entrepôt de données sur l'eau et le climat.} Il rassemble automatiquement quatre
sources publiques : les niveaux de nappe et les débits de rivière relevés par le réseau national,
la reconstitution du climat français depuis 1950, et le référentiel qui décrit les aquifères. Il
les met en cohérence, rattache chaque station de mesure au climat de son secteur, et se met à
jour chaque jour. Ce qui demandait plusieurs semaines de préparation à chaque équipe s'obtient
désormais par une simple requête.

\paragraph{Une plateforme d'exploration et de prévision.} Elle donne accès à trois usages. Un
\emph{observatoire} cartographique montre l'état des nappes, des rivières et du climat sur tout
le territoire. Un \emph{laboratoire de modélisation} permet de caler les modèles employés par les
hydrogéologues et de simuler des scénarios, par exemple l'effet d'un nouveau prélèvement
agricole. Un \emph{laboratoire d'intelligence artificielle} donne accès à une douzaine de
méthodes de prévision et, surtout, aux outils qui permettent d'examiner sur quoi leurs prévisions
reposent.

Ces réalisations répondent à trois attentes distinctes.

\begin{center}
\begin{tabularx}{\textwidth}{@{}lX@{}}
\toprule
\textbf{Aux hydrogéologues} & Des indicateurs de sécheresse normalisés, calculés de la même
manière sur tout le territoire, et un outil pour simuler l'effet d'un prélèvement sur une nappe
donnée. \\
\textbf{Aux chercheurs} & Un jeu de données prêt à l'emploi et un banc d'expérimentation, qui
suppriment le travail de préparation que chaque équipe refaisait pour son compte, avec ses
propres conventions. \\
\textbf{Au programme} & Une chaîne complète, de la donnée publique jusqu'au service explicable,
directement mobilisable pour les jumeaux numériques de la ressource en eau. \\
\bottomrule
\end{tabularx}
\end{center}

Trois résultats méritent d'être signalés.

\paragraph{Deux défauts de méthode ont été corrigés dans les indicateurs de sécheresse.} Le
premier faisait apparaître la France en manque d'eau permanent, y compris pendant les années
pluvieuses, parce qu'une grandeur climatique en remplaçait une autre qui porte un nom voisin mais
ne mesure pas la même chose. Le second empêchait de calculer l'indicateur sur un quart du
territoire, en raison d'un choix statistique inadapté à une partie des régions françaises. Sans
ces corrections, les cartes produites auraient été fausses, et l'erreur aurait été lue comme un
signal climatique. Le chapitre~\ref{ch:indicateurs} détaille les deux cas et les mesures qui les
établissent.

\paragraph{Une méthode d'explication qui ne propose que des situations possibles.} Expliquer une
prévision, c'est souvent répondre à « qu'aurait-il fallu pour que la nappe ne descende pas si
bas ». Les méthodes existantes répondent parfois par des combinaisons absurdes, comme un
réchauffement sans augmentation de l'évaporation. La méthode développée ici contraint la réponse
à rester physiquement cohérente, et à s'exprimer dans les termes du métier : \emph{un hiver vingt
pour cent plus sec et un degré plus chaud}. C'est la condition pour qu'une explication serve à
décider.

\paragraph{Une infrastructure éprouvée.} La chaîne traite plusieurs centaines de millions de
mesures, se relance sans rien dupliquer, conserve la trace de toute donnée écartée et reprend
après interruption. Ces propriétés déterminent si un jeu de données peut servir de référence
commune à plusieurs équipes, ou seulement à l'expérience qui l'a produit.

\medskip\noindent
L'ensemble constitue un socle réutilisable. Il est national, documenté, et son architecture n'est
pas propre aux eaux souterraines : le même enchaînement s'applique aux autres compartiments
environnementaux visés par le programme. Il est surtout destiné à servir de base aux travaux
d'explicabilité qui prendront la suite, auxquels il transmet à la fois les données, les modèles
et les outils d'analyse.

Deux réserves accompagnent cette conclusion. Aucune comparaison chiffrée des méthodes de
prévision n'est établie dans ce rapport : le temps consacré à l'infrastructure ne l'a pas permis,
et rien n'a été reconstitué de mémoire. Par ailleurs, l'usage de la plateforme reste aujourd'hui
circonscrit au cercle du projet ; son élargissement relève d'une décision d'accès plus que de
développements supplémentaires.
